\section{Conclusão e Trabalhos Futuros}

Este artigo apresentou uma avaliação experimental de servidores RTMP e transcodificadores em uma plataforma de streaming de vídeo ao vivo. A implementação permitiu comparar Nginx-RTMP e SRS combinados com FFmpeg e GStreamer sob cargas de 1, 10 e 50 viewers, com execução automatizada, coleta em CSV e agregação estatística de três repetições por tratamento.

Os resultados sustentam três conclusões principais dentro do ambiente estudado. Primeiro, o SRS disponibilizou a playlist mestre HLS mais cedo que o Nginx-RTMP: 8 s contra aproximadamente 22 s com FFmpeg e 1 s contra 22 s com GStreamer. Segundo, o GStreamer apresentou menor consumo de memória média, mas o Nginx + GStreamer apresentou maior pressão de CPU nas cargas de 10 e 50 viewers. Terceiro, SRS + GStreamer apresentou 3,15\% de erros HTTP no cenário de 50 viewers, acima do critério local de 3\%, além de elevada dispersão de CPU e memória. Assim, o menor tempo de criação da playlist não implica, por si só, maior estabilidade da sessão.

Essas conclusões descrevem as configurações implementadas e não devem ser generalizadas para qualquer hardware, codec ou parâmetro de codificação. O estudo foi executado em um único host, com codificação CPU-only, rede local, uma fonte sintética e alvos de bitrate não totalmente equivalentes entre FFmpeg e GStreamer. A latência ponta a ponta também não está validada: os campos disponíveis funcionam apenas como proxy de primeiro segmento. O resultado central deste estudo é, portanto, uma base experimental reproduzível e uma caracterização inicial de trade-offs, não um ranking universal das tecnologias.

Como continuidade, recomenda-se: normalizar os alvos de bitrate, parâmetros de encoder e número de representações; registrar versões e hashes das imagens; substituir o proxy por timestamps embutidos no conteúdo; coletar p50, p95 e p99; classificar os erros por recurso e código HTTP; testar cargas superiores; repetir a matriz em outros hosts; e avaliar aceleração por GPU. Essas etapas permitirão separar com maior segurança os efeitos do servidor RTMP, do transcodificador e da carga de espectadores.
